% =====================================================================
% 1 INTRODUÇÃO (prosa, sem subseções)
% Contém: contexto, teoria mínima (Fundamentação Teórica condensada, decisão 13),
% estado da arte e trabalhos correlatos (movidos da antiga seção 2 em 05/11/2026,
% regra do prof. Marlon: só 4 seções; modelo: TCC de Mateus S. de Matos, 2026-1), lacuna, problema e hipótese, objetivos (5 específicos, um verbo
% cada), justificativa e estrutura do trabalho.
% Origem de cada trecho: ver ALINHAMENTOS.md.
% =====================================================================

\section{INTRODUÇÃO}\label{sec:introducao}

O desempenho de sistemas de \textit{software} que armazenam dados em bancos relacionais depende da eficiência com que as consultas, escritas na linguagem SQL (\textit{Structured Query Language}), são executadas. Em aplicações corporativas, problemas de indexação, estimativas incorretas de cardinalidade, isto é, do número de registros que uma condição seleciona, e planos de execução lentos tornam-se críticos à medida que crescem o volume de dados e a complexidade das consultas \cite{sirigiri2025sqlperf}.

Consultas mal estruturadas, ausência de índices adequados ou uso incorreto dos índices existentes resultam em operações ineficientes, como as varreduras completas de tabela, que percorrem a tabela inteira mesmo quando apenas uma fração dos registros é necessária. Um número elevado de leituras lógicas, isto é, de páginas de dados lidas da memória durante a execução de uma consulta, indica essas varreduras e costuma decorrer da falta de índices \cite{rusanu2014perf}.

Para executar uma consulta, o SQL Server, sistema gerenciador de bancos de dados relacionais, conta com um componente interno, o otimizador de consultas, que escolhe a forma mais eficiente de acessar os dados. Em vez de executar a consulta exatamente como foi escrita, o otimizador a transforma em um plano de execução, uma sequência de operações físicas, como varreduras, buscas e junções (combinação de registros de tabelas diferentes), que descreve como os dados serão acessados \cite{fritchey2018plans}. A escolha do plano depende das estatísticas que o sistema mantém sobre a distribuição dos valores de cada coluna, as quais estimam quantos registros atendem a um filtro. Quando essas estimativas são imprecisas, por estatísticas desatualizadas ou por distribuições de dados irregulares, o otimizador pode escolher planos inadequados \cite{leis2015good, han2022cardest}.

O índice é uma estrutura auxiliar que organiza os valores de uma ou mais colunas para que o sistema localize registros sem percorrer a tabela inteira. Sua eficácia depende da seletividade da coluna indexada, a proporção entre valores distintos e registros: identificadores únicos têm alta seletividade e se beneficiam do índice, enquanto campos com poucos valores possíveis tendem a oferecer ganho reduzido \cite{fritchey2018qpt}. Em índices compostos, definidos sobre mais de uma coluna, a ordem das colunas determina o aproveitamento: se a consulta não filtra pela primeira coluna, o sistema recorre com frequência a uma varredura completa \cite{fritchey2018qpt}. O índice tem ainda um custo de manutenção, pois cada inserção, atualização ou exclusão precisa atualizá-lo, e o índice escolhido influencia o tempo de acesso e o de inserção \cite{saidu2024indexing}. Além dos índices, a forma como a consulta é escrita determina o plano resultante. Chamam-se antipadrões os padrões de escrita que retornam o resultado correto, mas levam a operações ineficientes \cite{karwin2010antipatterns}. Este trabalho tratou de quatro: a aplicação de funções sobre colunas usadas em filtros, que impede o uso eficiente do índice; a subconsulta correlacionada, reexecutada para cada registro da consulta externa e que pode ter tempo de execução quadrático em relação ao número de registros \cite{neumann2025unnesting}; a seleção desnecessária de colunas, que obriga o sistema a buscar na tabela base os dados que o índice não contém \cite{karwin2010antipatterns, fritchey2018qpt}; e a varredura completa por falta de índice, em que o sistema lê todas as páginas da tabela para encontrar poucos registros \cite{fritchey2018qpt}.

Apesar da relevância do tema, a escrita de consultas é uma tarefa exigente. \citeonline{dacosta2023sql} analisaram relatos de problema (\textit{issues}) de 20 projetos de código aberto da Apache Software Foundation e observaram que a mediana do tempo de resolução dos relatos relacionados a SQL foi de 25.897 minutos, contra 10.161 minutos nos demais, e que 53\% dos relatos relacionados a SQL receberam classificação de esforço percebido nas categorias mais altas da escala, contra 14\% nos demais, embora os próprios autores qualifiquem como pequena a diferença no tempo de resolução. Esse esforço favorece que a correção funcional, isto é, a capacidade de retornar os resultados esperados, prevaleça sobre o desempenho e que antipadrões como os descritos acima apareçam nas consultas.

Para situar este trabalho, foram analisadas pesquisas internacionais e nacionais sobre otimização de consultas e ajuste de índices, que se concentram na automação.

\citeonline{leis2021cost} propuseram um modelo que considera a carga de trabalho e o \textit{hardware} para determinar a configuração ideal de instâncias de computação em nuvem em sistemas de processamento analítico, observando que desempenho e custo podem variar em ordens de grandeza conforme a configuração escolhida. Sob a ótica deste trabalho, a limitação é que a decisão ocorre na camada de infraestrutura e não aborda a indexação nem a escrita da consulta em ambientes transacionais.

\citeonline{gretscher2025vldb} compararam abordagens de otimização de consultas, a dividida (\textit{split}), a holística e a híbrida, em um sistema de pesquisa, avaliando o tempo de otimização e a qualidade do plano de execução resultante, e propuseram uma estratégia híbrida. O estudo concentra-se na construção do plano pelo otimizador, ou seja, na perspectiva de quem desenvolve sistemas gerenciadores de bancos de dados, e não propõe método para quem escreve as consultas.

\citeonline{wu2024indextuning} elaboraram um levantamento sobre o ajuste automático de índices, que organiza as técnicas em pré-processamento, estimativa do benefício dos índices e seleção de índices, e discute sua aplicação em produtos comerciais e de código aberto. O levantamento é abrangente quanto às técnicas, mas não propõe um modelo estruturado que o desenvolvedor aplique de forma manual e iterativa ao escrever e revisar consultas.

\citeonline{brucato2024wred} propuseram a redução da carga de trabalho (\textit{workload reduction}) como técnica complementar para acelerar o ajuste de índices, diminuindo o tempo de cada chamada ao otimizador com índices hipotéticos sem afetar de modo relevante a qualidade do ajuste. Assim como o levantamento anterior, o trabalho mantém o foco na automação sobre uma carga de trabalho já existente, etapa posterior à escrita das consultas. \PENDENTE{confirmar com a orientadora se a autoria de pesquisadores da Microsoft em Brucato et al. (2024) é compatível com a regra de não citar a Microsoft (decisão 24)}

No contexto nacional, \citeonline{marins2021outertuning} apresentaram o \textit{Outer-Tuning}, \textit{framework} extensível, baseado em ontologia de domínio, para a sintonia semiautomática de bancos de dados relacionais. \citeonline{perciliano2021ondbtuning} detalharam a inferência de ações de sintonia com a ontologia OnDBTuning e a avaliaram na recomendação de índices. \citeonline{oliveira2022autonomic} propuseram um método automático para gerar e selecionar combinações de ações de sintonia, com soluções mais eficientes do que as locais independentes. Mais próximo da etapa de prescrição deste trabalho, \citeonline{araujo2014aresql} propuseram o ARe-SQL, abordagem automática e não intrusiva de reescrita de consultas SQL baseada em 11 heurísticas, com ganhos de até 92,5\% em alguns casos. A Tabela~\ref{tab:correlatos}, na seção~\ref{sec:comparacao}, resume os trabalhos e a lacuna de cada um.

Entre os métodos de ajuste automático de índices, parte considera apenas a recuperação de dados, sem o impacto das modificações, e outra parte inclui o custo de manutenção dos índices \cite{wu2024indextuning}. Nos planos coletados neste trabalho, a sugestão do recurso de índices ausentes do SQL Server traz as colunas e o impacto estimado na consulta, e não o custo que o índice impõe às operações de escrita.

A lacuna está em que, entre os trabalhos analisados, as soluções atuam sobre a consulta já escrita e visam à automação, e não foi encontrado método que oriente o desenvolvedor durante a escrita da consulta, relacione antipadrões a estratégias de indexação e verifique o efeito por métricas objetivas. O problema de pesquisa foi, portanto, como sistematizar boas práticas de indexação e de escrita de consultas em SQL Server de modo que seu efeito seja verificável. Formulou-se a hipótese de que a aplicação do \textit{framework} proposto reduz em pelo menos 20\% as leituras lógicas de consultas com antipadrões de escrita. \PENDENTE{confirmar com a orientadora o limiar de 20\%, que é parametrizável; o mesmo valor aparece no Resumo, no Abstract, nos objetivos, na seção 3 e na Conclusão} Para investigá-la, este trabalho propôs um \textit{framework} metodológico de quatro camadas sequenciais, diagnóstico, análise, prescrição e validação, fundamentado na literatura e validado em um experimento do tipo antes e depois com um conjunto de dados público.

O objetivo geral é propor um \textit{framework} metodológico de quatro camadas para a otimização de consultas em SQL Server, tendo como critério de eficácia a redução mediana de pelo menos 20\% nas leituras lógicas de consultas com antipadrões de escrita. Os objetivos específicos são revisar trabalhos correlatos qualificados sobre otimização de consultas e ajuste de índices, com as limitações de cada um consolidadas em quadro comparativo; catalogar quatro classes de antipadrões de escrita, cada uma com ao menos quatro consultas de teste nas versões original e otimizada; definir o modelo do \textit{framework}, com critérios verificáveis em cada camada e um limiar de melhoria significativa; executar o experimento antes e depois sobre o conjunto de dados público Stack Overflow 2010, com 20 repetições por consulta no cenário quente e 8 no frio; e mensurar a variação percentual das leituras lógicas, do tempo de processamento e do tempo de execução entre as versões original e otimizada, por antipadrão.

O trabalho se justifica pela contribuição de um método sistematizado, de aplicação manual e iterativa, que explicita a relação entre a escrita da consulta, a indexação e o desempenho, e que submete cada recomendação à verificação por métricas objetivas e à medição do custo de escrita. Entre os trabalhos analisados, o ineditismo está em reunir, em um único procedimento verificável, a identificação de antipadrões, a análise dos índices existentes, a prescrição e a validação. O SQL Server foi escolhido por oferecer recursos nativos de diagnóstico de desempenho, como o \textit{Query Store}, repositório que armazena o histórico dos planos de execução e permite identificar regressões de desempenho ao longo do tempo \cite{fritchey2018qpt}, e por adotar o T-SQL, dialeto da linguagem SQL próprio desse sistema.

Este artigo está organizado em quatro seções, contando esta introdução. A seção~\ref{sec:mm} descreve os materiais e métodos: o \textit{framework}, o catálogo de antipadrões, os dados e o ambiente, o delineamento, os procedimentos e a análise dos dados. A seção~\ref{sec:rd} apresenta e discute os resultados, e a seção~\ref{sec:conclusao} traz a conclusão, as contribuições, as limitações e os trabalhos futuros.
